% === COLORES PARA BORDES DE AGRUPACIONES ===
% Colores más intensos para bordes visibles
\definecolor{kgroup1}{RGB}{220,50,50}    % Rojo
\definecolor{kgroup2}{RGB}{50,180,50}    % Verde
\definecolor{kgroup3}{RGB}{50,50,220}    % Azul
\definecolor{kgroup4}{RGB}{200,150,0}    % Amarillo/Naranja
\definecolor{kgroup5}{RGB}{180,50,180}   % Magenta
\definecolor{kgroup6}{RGB}{0,180,180}    % Cian

% === MACRO PARA DIBUJAR UN GRUPO CON BORDE ===
% Uso: \karnaughGrupo{fila_inicio}{col_inicio}{filas}{columnas}{color}
% Las coordenadas son 0-indexed (fila 0-3, col 0-3)
\newcommand{\karnaughGrupo}[5]{%
  % #1 = fila inicio, #2 = col inicio, #3 = num filas, #4 = num cols, #5 = color
  \draw[#5, line width=2.5pt, rounded corners=6pt]
    (#2*\cellsize + 0.08, -#1*\cellsize - 0.08) 
    rectangle 
    ++(#4*\cellsize - 0.16, -#3*\cellsize + 0.16);
}

% === MACRO PRINCIPAL PARA KARNAUGH DE 4 VARIABLES CON GRUPOS ===
% Uso: \karnaughConGrupos{AB}{CD}{v0,v1,...,v15}{código TikZ para grupos}
% El cuarto parámetro permite dibujar grupos con \karnaughGrupo
\newcommand{\karnaughConGrupos}[4]{%
\begin{tikzpicture}[scale=0.8]
  \def\cellsize{1.2}
  
  % Grid principal (fondo blanco)
  \fill[white] (0,0) rectangle (4*\cellsize,-4*\cellsize);
  
  % Líneas del grid
  \draw[thick] (0,0) rectangle (4*\cellsize,-4*\cellsize);
  \draw (0,-\cellsize) -- (4*\cellsize,-\cellsize);
  \draw (0,-2*\cellsize) -- (4*\cellsize,-2*\cellsize);
  \draw (0,-3*\cellsize) -- (4*\cellsize,-3*\cellsize);
  \draw (\cellsize,0) -- (\cellsize,-4*\cellsize);
  \draw (2*\cellsize,0) -- (2*\cellsize,-4*\cellsize);
  \draw (3*\cellsize,0) -- (3*\cellsize,-4*\cellsize);
  
  % Etiqueta de variables (filas)
  \node[font=\bfseries] at (-0.6, -2*\cellsize) {#1};
  % Etiqueta de variables (columnas)  
  \node[font=\bfseries] at (2*\cellsize, 0.6) {#2};
  
  % Código Gray para columnas (CD)
  \node at (0.5*\cellsize, 0.25) {00};
  \node at (1.5*\cellsize, 0.25) {01};
  \node at (2.5*\cellsize, 0.25) {11};
  \node at (3.5*\cellsize, 0.25) {10};
  
  % Código Gray para filas (AB)
  \node at (-0.35, -0.5*\cellsize) {00};
  \node at (-0.35, -1.5*\cellsize) {01};
  \node at (-0.35, -2.5*\cellsize) {11};
  \node at (-0.35, -3.5*\cellsize) {10};
  
  % Colocar valores en las celdas
  \foreach \val [count=\i from 0] in {#3} {
    \pgfmathtruncatemacro{\row}{int(\i/4)}
    \pgfmathtruncatemacro{\col}{mod(\i,4)}
    \node at (\col*\cellsize + 0.5*\cellsize, -\row*\cellsize - 0.5*\cellsize) {\val};
  }
  
  % Dibujar los grupos (bordes coloreados) ENCIMA del grid
  #4
\end{tikzpicture}%
}

% === MACRO SIMPLIFICADA SIN GRUPOS ===
\newcommand{\karnaughSimple}[3]{%
  \karnaughConGrupos{#1}{#2}{#3}{}%
}

% === MACRO LEGACY (compatibilidad con el formato anterior) ===
% Convierte el formato antiguo fila/col/color a bordes
% NOTA: Esta macro mantiene compatibilidad pero usa colores de fondo
\newcommand{\karnaughCuatro}[4]{%
\begin{tikzpicture}[scale=0.8]
  \def\cellsize{1.2}
  
  % Dibujar celdas con colores de fondo (formato legacy)
  \foreach \row/\col/\bgcolor in {#4} {
    \fill[\bgcolor, opacity=0.3] (\col*\cellsize, -\row*\cellsize) rectangle ++(\cellsize,-\cellsize);
    \draw[\bgcolor, line width=2pt, rounded corners=4pt] 
      (\col*\cellsize + 0.06, -\row*\cellsize - 0.06) 
      rectangle ++(\cellsize - 0.12, -\cellsize + 0.12);
  }
  
  % Grid principal
  \draw[thick] (0,0) rectangle (4*\cellsize,-4*\cellsize);
  \draw (0,-\cellsize) -- (4*\cellsize,-\cellsize);
  \draw (0,-2*\cellsize) -- (4*\cellsize,-2*\cellsize);
  \draw (0,-3*\cellsize) -- (4*\cellsize,-3*\cellsize);
  \draw (\cellsize,0) -- (\cellsize,-4*\cellsize);
  \draw (2*\cellsize,0) -- (2*\cellsize,-4*\cellsize);
  \draw (3*\cellsize,0) -- (3*\cellsize,-4*\cellsize);
  
  % Etiqueta de variables
  \node[font=\bfseries] at (-0.6, -2*\cellsize) {#1};
  \node[font=\bfseries] at (2*\cellsize, 0.6) {#2};
  
  % Código Gray para columnas
  \node at (0.5*\cellsize, 0.25) {00};
  \node at (1.5*\cellsize, 0.25) {01};
  \node at (2.5*\cellsize, 0.25) {11};
  \node at (3.5*\cellsize, 0.25) {10};
  
  % Código Gray para filas
  \node at (-0.35, -0.5*\cellsize) {00};
  \node at (-0.35, -1.5*\cellsize) {01};
  \node at (-0.35, -2.5*\cellsize) {11};
  \node at (-0.35, -3.5*\cellsize) {10};
  
  % Colocar valores
  \foreach \val [count=\i from 0] in {#3} {
    \pgfmathtruncatemacro{\row}{int(\i/4)}
    \pgfmathtruncatemacro{\col}{mod(\i,4)}
    \node at (\col*\cellsize + 0.5*\cellsize, -\row*\cellsize - 0.5*\cellsize) {\val};
  }
\end{tikzpicture}%
}
